\section{Why smooth optimizers fail there}
\label{sec:smooth}

L-BFGS builds a curvature model from successive gradient differences and relies on
a line search that brackets a point satisfying the Wolfe conditions
\citep{nocedal2006numerical}, whose convergence theory for BFGS assumes a twice
continuously differentiable objective \citep[Assumption~6.1]{nocedal2006numerical}.
Both assume a gradient that varies continuously. Across a kink the gradient
differences carry no information about curvature, and the bracketing interval can
collapse without a sufficient-decrease point inside it.

BFGS is not helpless on nonsmooth functions. \citet{lewis2013nonsmooth} find that,
with a line search built on the \emph{weak} Wolfe condition, it consistently
converges to local minimisers in their experiments, at a linear rate. They are as
clear that ``for nonsmooth optimization, it is clear that enforcing the strong Wolfe
condition is not possible in general'', and suggest that L-BFGS codes offer a
weak-Wolfe line search for that reason. Stan's line search imposes the strong
conditions (\code{stan/optimization/bfgs_linesearch.hpp}), and Stan's L-BFGS differs
from their method in its limited memory and in stopping tests on changes in the
objective, the gradient and the parameters as well. We do not try to separate
which of these matters most; \cref{sec:evidence} establishes what the combination
does. The failure is quiet: in every fit of \cref{fig:margin} where Prophet ran
L-BFGS it terminated normally, without the fall-back to Newton that Prophet
triggers on an abnormal exit.

Newton's method is exposed for a different reason. It needs the Hessian, and
along $\{\delta_j = 0\}$ the curvature of $\lvert\delta_j\rvert$ is not small but
undefined. A second-order method that also stops short of the optimum is evidence
that the obstacle is the objective rather than any one algorithm, which is why
\cref{sec:evidence} reports both.

Methods built for this structure exist. Orthant-wise limited-memory quasi-Newton
\citep[OWL-QN;][]{andrew2007scalable} handles an $\ell_1$ term directly by
restricting each step to an orthant. The approach taken in \cref{sec:reformulation}
is older and simpler: remove the non-smoothness from the objective rather than
from the method.
